% ============================================================
% Methodology -- DRAFT NOTES
% Not a final thesis chapter. Paste into Overleaf as-is.
% No em-dashes used anywhere in this document.
% ============================================================

\documentclass[12pt, a4paper]{article}

\usepackage[margin=2.5cm]{geometry}
\usepackage{booktabs}
\usepackage{parskip}
\usepackage{enumitem}
\usepackage{xcolor}
\usepackage{fancyhdr}
\usepackage{array}
\usepackage{hyperref}

\pagestyle{fancy}
\fancyhf{}
\rhead{\textcolor{gray}{\small DRAFT -- Methodology}}
\lhead{\textcolor{gray}{\small JS/TS Coding Agent Adoption Study}}
\cfoot{\thepage}

\title{%
  \textbf{Methodology Draft}\\[0.4em]
  \large Measuring AI Coding Agent Adoption Across Countries\\
  in the JavaScript/TypeScript/npm Ecosystem on GitHub
}
\author{Manuel Youssef Haik Kevorkian}
\date{October 2026}

\begin{document}
\maketitle
\tableofcontents
\newpage

% ============================================================
\section{Overview and Research Questions}
% ============================================================

This study measures the adoption of AI coding agents across GitHub repositories
in the JavaScript and TypeScript ecosystem, with the goal of ranking countries
by their adoption rates and testing whether differences across countries are
statistically significant. The study is designed as an observational empirical
study: it does not control who uses AI agents or intervene in any way, but
instead detects adoption from publicly observable signals left in repository
histories.

\textbf{Research Question 1 (Global adoption).}
What is the overall AI coding agent adoption rate in the JavaScript/TypeScript/npm
ecosystem on GitHub, and is it statistically non-trivial?

\textbf{Research Question 2 (Undercounting).}
How much does single-signal detection undercount adoption compared to a
multi-signal approach, and is this difference statistically significant?

\textbf{Research Question 3 (Geographic variation).}
Do adoption rates differ significantly across countries, and which countries
show the highest and lowest rates within the JS/TS/npm ecosystem?

These questions build directly on the gaps identified in the literature.
Robbes et al.\ found a 6.46 pp undercounting delta across all GitHub languages,
and Khosravani and Mockus documented a 30$\times$ gap from single-signal
approaches. Neither study measured the JS/TS ecosystem specifically.
Daniotti et al.\ showed country-level variation in AI tool adoption but focused
on Python and function-level analysis, not multi-signal repository mining.

% ============================================================
\section{Study Design and Scope}
% ============================================================

\subsection{Unit of Analysis}

The unit of analysis is the individual GitHub repository. A repository is
the natural unit for this study because it is the level at which developers
configure AI agents (via configuration files at the root), commit code (where
trailer patterns appear), and open pull requests (where PR branch names and
labels are visible). Country-level adoption rates are then computed by
aggregating repository-level labels within each country.

\subsection{Scope Boundaries}

\begin{itemize}
  \item \textbf{Language ecosystem:} JavaScript and TypeScript repositories
        with at least 175 KB of combined JS/TS code, confirmed to be npm
        projects by the presence of \texttt{package.json} alongside a lockfile.
  \item \textbf{Time scope:} Only repositories created on or after 24 February
        2025 are included. This date marks the point at which AI coding agent
        tooling reached broad mainstream availability in the npm ecosystem,
        with Claude Code launching publicly and Cursor, Windsurf, and GitHub
        Copilot Workspace all in active use.
  \item \textbf{Platform:} Public repositories on GitHub only. Private
        repositories, repositories on other platforms, and local codebases
        are out of scope.
  \item \textbf{Detection:} Only repositories that leave an observable
        artefact (configuration file, commit trailer, PR marker, or SDK
        dependency) are counted as adopted. Silent use of AI tools is not
        captured. This means the study reports a lower bound on true adoption.
\end{itemize}

% ============================================================
\section{Phase 1: Corpus Construction}
% ============================================================

\subsection{Repository Collection}

Repositories are collected using the GitHub GraphQL Search API. Because the
API caps results at 1,000 per query, the search is partitioned into 5-day
creation date windows covering the full study period. Any window returning
more than 1,000 results is split recursively in half until every window fits
within the cap. This approach, following the methodology of Dabi\'c et
al.~\cite{dabic2021}, ensures complete enumeration without missing
repositories.

The search is run twice: once for the \texttt{language:JavaScript} filter and
once for \texttt{language:TypeScript}, with results deduplicated on
\texttt{full\_name} (owner/repo). This captures repositories whose primary
language could be classified as either by GitHub's language detection.

\subsection{Inclusion Filters}

The following filters are applied at query time and reflect criteria agreed
across the research team for the broader multi-ecosystem study:

\begin{table}[h]
\centering
\caption{Agreed corpus inclusion filters.}
\begin{tabular}{@{}lll@{}}
\toprule
\textbf{Filter} & \textbf{Value} & \textbf{Rationale} \\
\midrule
Stars           & $\geq 4$              & Remove placeholders and abandoned repos \\
Commits         & $\geq 50$             & Ensure meaningful development activity \\
Fork            & \texttt{false}        & Original repositories only \\
Creation date   & $\geq$ 2025-02-24     & Repositories born in the AI agent era \\
JS+TS bytes     & $\geq 175\,000$ bytes & Language must be primary, not incidental \\
\bottomrule
\end{tabular}
\end{table}

The star and commit thresholds are consistent with the quality filters
recommended by Kalliamvakou et al.~\cite{kalliamvakou2016}, who found that
single-committer, short-lived, and low-activity repositories constitute the
majority of GitHub but do not represent active software development.

\subsection{Ecosystem Categorization}

After collection, each repository is categorized by its package manager.
A single GraphQL batch query per 15 repositories checks for the presence of
seven root files: \texttt{package.json}, \texttt{package-lock.json},
\texttt{yarn.lock}, \texttt{pnpm-lock.yaml}, \texttt{bun.lockb},
\texttt{deno.json}, and \texttt{deno.jsonc}. No file content is read; only
existence is checked.

\begin{table}[h]
\centering
\caption{Ecosystem categories assigned in Phase 1.}
\begin{tabular}{@{}ll@{}}
\toprule
\textbf{Category} & \textbf{Detection criterion} \\
\midrule
npm          & \texttt{package.json} + \texttt{package-lock.json} \\
yarn         & \texttt{package.json} + \texttt{yarn.lock} \\
pnpm         & \texttt{package.json} + \texttt{pnpm-lock.yaml} \\
bun          & \texttt{package.json} + \texttt{bun.lockb} \\
npm\_no\_lock & \texttt{package.json} present, no lockfile found \\
deno         & \texttt{deno.json} or \texttt{deno.jsonc} present \\
vanilla      & None of the above files found \\
\bottomrule
\end{tabular}
\end{table}

Priority order when multiple lockfiles are present: npm $>$ yarn $>$ pnpm $>$ bun.
The \texttt{ecosystem\_type} column is added to the corpus for downstream
filtering if needed.

\subsection{Output}

\texttt{data/stage1/global\_corpus.csv}: all repositories passing the filters,
with fields \texttt{full\_name}, \texttt{owner\_login}, \texttt{repo\_name},
\texttt{created\_at}, \texttt{stars}, \texttt{commits}, \texttt{lang\_bytes},
\texttt{primary\_language}, \texttt{size\_stratum}, \texttt{ecosystem\_type}.

% ============================================================
\section{Phase 2: Owner Location Inference}
% ============================================================

\subsection{Motivation}

GitHub does not expose a location field at the repository level. Location is
attached to user profiles. For each unique owner login in the corpus, a
country is inferred in two passes. This approach is validated by Wachs et
al.~\cite{wachs2022}, who confirmed that self-reported GitHub profile location
is broadly consistent with other geographic signals for developers who fill
in the field.

\subsection{Pass 1: Fast Location Lookup}

The GitHub GraphQL API is queried in batches of 50 owners per request to
fetch the raw \texttt{location} profile field. Each raw string is matched
against a keyword lookup table covering country names, major cities, and
common abbreviations in multiple languages (English and Arabic included).
Matching uses substring search with the longest key checked first to prevent
partial matches (e.g., ``New Zealand'' is checked before ``New'').

Owners whose field is empty, contains only generic terms (``Earth'',
``Remote''), or does not match any keyword are marked unresolved and passed
to Pass 2.

\subsection{Pass 2: Deep Profile Inference}

For each unresolved owner, one GraphQL query fetches additional profile
fields. Six signals are tried in the order shown below; inference stops at
the first successful match:

\begin{table}[h]
\centering
\caption{Signal cascade used in Pass 2 of location inference.}
\begin{tabular}{@{}clll@{}}
\toprule
\textbf{Priority} & \textbf{Signal} & \textbf{Profile field} & \textbf{Confidence} \\
\midrule
1 & Bio keyword match      & \texttt{bio}                      & Medium \\
2 & Company keyword match  & \texttt{company}                  & Medium \\
3 & Company domain TLD     & \texttt{company} (domain)         & Medium \\
4 & Email domain TLD       & \texttt{email}                    & Medium \\
5 & Website URL TLD        & \texttt{websiteUrl}               & Low    \\
6 & Org location keyword   & \texttt{organizations.location}   & Medium \\
7 & Network majority       & \texttt{followers} + \texttt{following} & Low \\
\bottomrule
\end{tabular}
\end{table}

Country-code top-level domains (.de for Germany, .eg for Egypt, .br for
Brazil, and so on) are used for signals 3, 4, and 5. Network majority
requires at least 3 located connections with at least 40\% agreeing on the
same country. This threshold prevents a small number of connected accounts
from dominating the inference.

Owners that remain unresolved after all seven signals are labelled
\texttt{not\_specified} and excluded from the country-level analysis in
Phase 4. No country is forced or guessed beyond these signals, consistent
with the approach taken by Wachs et al.~\cite{wachs2022}.

\subsection{Output}

\texttt{data/stage2/owner\_locations\_all.csv}: every unique owner login with
fields \texttt{inferred\_country}, \texttt{confidence}, and \texttt{source}.
This file is the master input for Phase 4.

% ============================================================
\section{Phase 3: AI Adoption Detection}
% ============================================================

\subsection{Design Principle}

The detection pipeline uses three passes in ascending order of API cost.
A repository classified by an earlier pass is never re-examined by later
passes. This ensures that the most expensive checks run only on the smallest
possible subset of repositories. The design follows the layered approach
recommended by Robbes et al.~\cite{robbesHeuristics2026} and Mao et
al.~\cite{mao2026}, who showed that cheap signals (configuration files)
classify the majority of positive cases instantly.

This thesis uses \emph{heuristic-based} detection rather than
machine-learning classification, for three reasons. First, no labelled
training dataset exists for the JS/TS/npm ecosystem. Second, heuristic
approaches are fully transparent and reproducible. Third, they are directly
comparable to the methods used by Robbes et al.\ and Khosravani and
Mockus~\cite{khosravani2026}, which are the primary baseline studies. The
ML-based approach of Ghaleb~\cite{ghaleb2026} is noted as a limitation:
the heuristic approach will miss repositories where agents are used but no
observable artefact is left behind.

\subsection{Pass 1: Configuration File Check}

Every repository is checked for the presence of 12 known AI agent
configuration files at the repository root. Checks are batched (10
repositories per GraphQL query), making this pass extremely efficient: fewer
than 2,400 API calls cover all 60,129 repositories. These files are
\emph{hard positives} in the terminology of Chatlatanagulchai et
al.~\cite{chatlatanagulchai2025}: no developer accidentally creates a
\texttt{CLAUDE.md} or \texttt{.cursorrules} file.

\begin{table}[h]
\centering
\caption{Configuration files checked in Pass 1.}
\begin{tabular}{@{}ll@{}}
\toprule
\textbf{File} & \textbf{Agent} \\
\midrule
\texttt{CLAUDE.md}                        & Claude Code \\
\texttt{AGENTS.md}                        & Claude Code / Generic \\
\texttt{.cursorrules}                     & Cursor \\
\texttt{.cursorignore}                    & Cursor \\
\texttt{.aider.conf.yml}                  & Aider \\
\texttt{.roomodes}                        & Roo Code \\
\texttt{.windsurfrules}                   & Windsurf \\
\texttt{GEMINI.md}                        & Gemini CLI \\
\texttt{.github/copilot-instructions.md}  & GitHub Copilot \\
\texttt{.continue/config.json}            & Continue \\
\texttt{.codeium/config.json}             & Codeium \\
\texttt{AGENTS.json}                      & OpenAI Codex \\
\bottomrule
\end{tabular}
\end{table}

Any match immediately labels the repository \textbf{agentic}. No further
passes are needed.

\subsection{Pass 2: Commit Trailer and PR Branch Scan}

Repositories not classified in Pass 1 have their 100 most recent commits
fetched and scanned for commit message trailer patterns. The 50 most recent
pull requests are also checked for agent-specific branch name prefixes. This
approach is based directly on the heuristic catalogue of Robbes et
al.~\cite{robbesHeuristics2026}, extended with additional patterns for agents
that became prominent after that study.

The bot account detection approach of Dey et al.~\cite{dey2020biman} also
informs Pass 2: registered GitHub app accounts (e.g.\
\texttt{github-copilot[bot]}, \texttt{devin-ai-integration[bot]}) are
checked via author email patterns alongside the commit message trailers.

Representative patterns are listed below. The full list is in the source code.

\begin{table}[h]
\centering
\caption{Representative commit and PR patterns in Pass 2.}
\begin{tabular}{@{}ll@{}}
\toprule
\textbf{Pattern} & \textbf{Agent} \\
\midrule
\texttt{Co-authored-by: ...claude...anthropic} & Claude Code \\
\texttt{noreply@anthropic.com} (author email)  & Claude Code \\
\texttt{Co-authored-by: ...github.copilot}     & GitHub Copilot \\
\texttt{Co-authored-by: ...devin}              & Devin \\
\texttt{Co-authored-by: ...cursor}             & Cursor \\
\texttt{Co-authored-by: ...aider}              & Aider \\
\texttt{Co-authored-by: ...windsurf}           & Windsurf \\
\texttt{Co-authored-by: ...gemini}             & Gemini \\
PR branch prefix \texttt{copilot/}            & GitHub Copilot \\
PR branch prefix \texttt{codex/}              & OpenAI Codex \\
PR branch prefix \texttt{devin/}              & Devin \\
\bottomrule
\end{tabular}
\end{table}

\subsection{Pass 3: Deep Scan}

Repositories not classified by Pass 1 or Pass 2 receive a comprehensive
scan. This pass runs on the smallest subset and applies the most thorough
checks: up to 300 commits (message and author email), up to 100 pull
requests (body text, labels, and branch name), and the
\texttt{package.json} dependency file (checking for AI agent SDK
dependencies such as \texttt{@anthropic-ai/sdk} and \texttt{@cursor/agent}).

The npm dependency signal is an additional channel not available in other
language ecosystems. It detects repositories that integrate AI agent SDKs as
software dependencies rather than as development tools, which represents a
distinct category of adoption not captured by commit trailers or configuration
files alone.

\subsection{Classification Rule}

The same classification rule applies across all three passes:

\begin{table}[h]
\centering
\caption{Adoption classification thresholds.}
\begin{tabular}{@{}ll@{}}
\toprule
\textbf{Condition} & \textbf{Label} \\
\midrule
$\geq 3$ AI signals, OR signals from $\geq 2$ distinct agents & \textbf{agentic} \\
1 or 2 AI signals from a single agent                         & \textbf{ai\_assisted} \\
No signal found across all three passes                        & \textbf{none} \\
\bottomrule
\end{tabular}
\end{table}

This two-level classification follows the framework of Galster et
al.~\cite{galsterConfig2026} and Chatlatanagulchai et
al.~\cite{chatlatanagulchai2025}: repositories with isolated signals may
reflect occasional use, while repositories with multiple signal types or
agents reflect sustained integration.

\subsection{Multi-Signal vs Single-Signal Comparison (RQ2)}

To answer RQ2, the study records the label that would have been assigned
using Pass 1 alone (configuration files only) alongside the label assigned
by the full three-pass pipeline. The gap between these two classification
schemes, measured via McNemar's test on the paired per-repository outcomes,
gives the undercounting delta for the JS/TS/npm ecosystem.

\subsection{Output}

\texttt{data/stage3/corpus\_detected\_full.csv}: the full corpus with fields
\texttt{adoption\_label} (agentic / ai\_assisted / none),
\texttt{agents\_found}, and \texttt{signals\_found} for every repository.

% ============================================================
\section{Phase 4: Geographic Analysis}
% ============================================================

\subsection{Joining Repos with Countries}

The detected corpus is joined with \texttt{owner\_locations\_all.csv} on
\texttt{owner\_login}. Each repository inherits its owner's inferred country.
Repositories whose owner has no country label (\texttt{not\_specified}) are
retained in global totals but excluded from country-level analysis.

\subsection{Country-Level Aggregation}

For each country with at least 30 repositories in the corpus, the following
metrics are computed: total repositories, adopted count (agentic +
ai\_assisted), agentic-only count, adoption rate (\%), agentic rate (\%),
and the 95\% Wilson confidence interval for the adoption rate. Countries are
ranked by adoption rate in descending order to produce the primary result
table. The 30-repository threshold balances the desire to include as many
countries as possible against the need for meaningful confidence intervals.

\subsection{Statistical Tests}

Three formal statistical tests are performed, one per research question:

\textbf{RQ1 (global adoption).} A one-sample binomial test against the null
hypothesis that the global adoption rate is at most 5\%. A 5\% null rate
corresponds to the level one would expect from incidental tool use (e.g.\
a developer trying an agent once and never again). Rejecting this null
confirms that adoption is non-trivial in absolute terms.

\textbf{RQ2 (undercounting).} McNemar's test on paired per-repository
detection outcomes (Pass 1 only vs.\ full three-pass pipeline). McNemar's
test is appropriate here because the two classification schemes are applied
to the same repositories, making the outcomes paired. The test statistic uses
the continuity correction, and the undercounting delta is reported as the
difference in percentages.

\textbf{RQ3 (geographic variation).} A chi-squared test of independence
across all eligible countries tests whether adoption rates differ
significantly across the country groups. Because the chi-squared test does
not identify which specific countries differ, it is followed by pairwise
Fisher's exact tests comparing each country to the global baseline. Multiple
comparisons are corrected using the Benjamini-Hochberg false-discovery rate
procedure at $\alpha = 0.05$.

% ============================================================
\section{Validity Threats}
% ============================================================

\textbf{Construct validity.} The study measures adoption through observable
artefacts, not through developer surveys or interviews. Developers who use
agents without leaving any of the detectable signals are not counted. The
multi-signal approach mitigates this by expanding the set of detectable
patterns, but silent use remains invisible. The reported rates are therefore
lower bounds on true adoption.

\textbf{Internal validity.} The geographic inference for owner location
relies on self-reported profile data, which may be inaccurate, out of date,
or absent. Approximately 27\% of owners are resolved by the fast pass;
the deep pass recovers additional owners but an estimated 30--40\% of owners
remain unresolved. If unresolved owners differ systematically from resolved
ones (for example, if developers from lower-income countries are less likely
to fill in their profiles, as Wachs et al.~\cite{wachs2022} suggest), the
country-level rates will be biased. This limitation is reported explicitly
alongside results.

\textbf{External validity.} The corpus is restricted to the
JavaScript/TypeScript/npm ecosystem on GitHub. Results may not generalise to
other language ecosystems, to private repositories, or to platforms other
than GitHub.

\textbf{Reliability.} All detection patterns are deterministic regex matches
applied to publicly accessible repository data. The full pipeline is
open-source and reproducible: given the same GITHUB\_TOKEN and the same date
range, the same corpus and the same adoption labels should be produced.

\end{document}
